% 正文片段。移入主论文时需加载 amsmath、amssymb、booktabs、tikz、float，
% 并加载 TikZ arrows.meta 库；可将下面的分页宏设为空以随全文自然排版。
\providecommand{\coveragepagebreak}{}
\subsection*{定位区域的同直径圆覆盖分析}

给定前文求得的非空有界定位区域 $P$ 的直径 $D$ 及构成直径的最远顶点 $A,B$。以下讨论是否存在直径为 $D$ 的圆覆
盖整个区域。值得注意的是，当区域退化为单点时，半径为零的圆盘即可覆盖；
下文仅考虑 $D>0$ 的情形。

\subsubsection*{圆心确定与顶点覆盖判据}

设某个半径为 $D/2$、圆心为 $O$ 的圆盘能够覆盖 $P$，则它必须同时包含 $A,B$。由三角不等式，
\begin{equation}\label{eq:cov-center}
D=\|A-B\|_2\le\|A-O\|_2+\|O-B\|_2\le D.
\end{equation}
两端相等，要求 $O$ 在线段 $AB$ 上，且到两端点的距离均为 $D/2$，因此圆心只能取
\begin{equation}
C=\frac{A+B}{2}.
\end{equation}
这说明，同直径圆能否覆盖区域并不是任意选取圆心的问题：若以 $C$ 为圆心的圆盘不能覆盖，则移动圆心也无法在保持直径 $D$ 的同时实现覆盖。

设区域顶点为 $V_0,\ldots,V_{h-1}$，定义
\begin{equation}
\rho=\max_{0\le i<h}\|V_i-C\|_2,\qquad \Delta=\rho-D/2.
\end{equation}
圆盘是凸集，包含全部顶点就包含它们的凸包，即整个定位区域。因此，
\begin{equation}\label{eq:cov-check}
\boxed{\text{存在同直径圆盘覆盖 }P\quad\Longleftrightarrow\quad\rho\le D/2.}
\end{equation}
由于 $A,B$ 已在候选圆周上，精确运算中 $\rho\ge D/2$，故覆盖时 $\Delta=0$；若 $\Delta>0$，达到 $\rho$ 的顶点就是不能覆盖的直接证据。已知直径端点后，只需遍历全部顶点，检验复杂度为 $O(h)$。

\subsubsection*{由测向角域构造反例}

为判断“是否总能覆盖”，构造三次测向形成等边三角形的反例。设干扰源真实位置为原点，检测点及示向度如表\ref{tab:cov-input}所示，各测向角域半宽均为 $1^\circ$。
\begin{table}[H]
\centering
\caption{等边三角形反例的测向输入（坐标单位：米）}\label{tab:cov-input}
\begin{tabular}{crrr}
\toprule
检测点 & $x_i$ & $y_i$ & 示向度/$(^\circ)$\\
\midrule
$S_1$ & $500$ & $500\sqrt3+20/\sqrt3$ & $241$\\
$S_2$ & $-1010$ & $-10/\sqrt3$ & $1$\\
$S_3$ & $510$ & $-510\sqrt3+20/\sqrt3$ & $121$\\
\bottomrule
\end{tabular}
\end{table}
这些角域的下边界分别与三角形的三条边共线，上边界不进一步切割三角形，其交集恰为
\begin{equation}
A=\left(-10,-\frac{10}{\sqrt3}\right),\quad
B=\left(10,-\frac{10}{\sqrt3}\right),\quad
T=\left(0,\frac{20}{\sqrt3}\right)
\end{equation}
所围成的等边三角形。原点位于三个角域内，各检测点距源约为 $1000$ 米，选择 $1200$ 米有效接收半径即可满足题设接收条件；源与定位三角形亦均位于目标圆域内。因此，该反例可由题设交会定位过程实际构造，并非任意指定的多边形。

\coveragepagebreak

该三角形边长为 $20$ 米，故 $D=20$ 米。取底边 $AB$ 为直径端点，候选圆的圆心为 $C=(0,-10/\sqrt3)$、半径为 $10$ 米，而
\begin{equation}
\|T-C\|_2=10\sqrt3\approx17.3205>10.
\end{equation}
顶点 $T$ 超出圆盘约 $7.3205$ 米，如图\ref{fig:cov-counterexample}所示。结合圆心的唯一性，任意直径为 $20$ 米的圆盘均无法覆盖该定位区域。

\begin{figure}[H]
\centering
% 矢量几何图；坐标单位按米等比例绘制，缩放不损失清晰度。
\begin{tikzpicture}[x=0.20cm,y=0.20cm,>=Latex,font=\small]
\definecolor{covblue}{RGB}{43,85,133}
\definecolor{covred}{RGB}{191,54,55}
\definecolor{covgreen}{RGB}{20,125,100}
\coordinate (A) at (-10,{-10/sqrt(3)});
\coordinate (B) at (10,{-10/sqrt(3)});
\coordinate (T) at (0,{20/sqrt(3)});
\coordinate (C) at (0,{-10/sqrt(3)});
\coordinate (O) at (0,0);
\fill[covblue!9] (A)--(B)--(T)--cycle;
\draw[covgreen,line width=1.1pt] (O) circle[radius={20/sqrt(3)}];
\draw[covred,dashed,line width=1.1pt] (C) circle[radius=10];
\draw[covblue,line width=1pt] (A)--(B)--(T)--cycle;
\draw[covred,line width=1.3pt] (A)--(B);
\draw[covblue,densely dotted] (C)--(T);
\foreach \p in {A,B,T}{\fill[covblue] (\p) circle[radius=0.22];}
\fill[covred] (C) circle[radius=0.2];
\fill[covgreen] (O) circle[radius=0.2];
\node[below left] at (A) {$A$};
\node[below right] at (B) {$B$};
\node[above] at (T) {$T$（圆外顶点）};
\node[below right] at (C) {$C$};
\node[left] at (O) {$O_*$};
\node[above,fill=white,inner sep=1pt] at (-5,{-10/sqrt(3)}) {$D=20$};
\draw[<->,thin] (12.5,{-10/sqrt(3)})--(12.5,{20/sqrt(3)}) node[midway,right] {$CT=10\sqrt3>10$};
\draw[covred,dashed,line width=1.1pt] (19,5)--(23,5);
\node[anchor=west] at (24,5) {同直径圆：$r=10$};
\draw[covgreen,line width=1.1pt] (19,0)--(23,0);
\node[anchor=west] at (24,0) {最小包围圆：$R_*=20/\sqrt3$};
\node[anchor=west,text=covblue] at (19,-6) {蓝色区域：可行定位区域};
\end{tikzpicture}

\caption{等边三角形反例：红色虚线圆盘不能覆盖顶点 $T$，绿色实线为最小包围圆。}\label{fig:cov-counterexample}
\end{figure}

\subsubsection*{最小包围圆解释与数值核验}

令 $R_*=\min_O\max_i\|V_i-O\|_2$ 为最小包围圆半径。任何覆盖圆均需包含最远点对，故 $R_*\ge D/2$，同直径覆盖成立当且仅当 $R_*=D/2$。有限点集的最小包围圆由至多三个边界支撑点确定\textsuperscript{[1]}：可由两点构成直径确定，或由三个不可再约简的支撑点构成的锐角三角形之外接圆确定。直角三角形可归入斜边两端支撑的情形；圆周上的点数可以多于三个。

程序据此枚举两点直径圆及三点外接圆，检验其是否包含全部顶点后取最小半径，用于独立核验覆盖判据。对本反例，$R_*=20/\sqrt3>10$，最小覆盖圆直径约为 $23.0940$ 米。代码中另有两次测向形成的四边形反例及可覆盖对照例，结果如下。
\begin{table}[H]
\centering
\caption{区域直径与最小覆盖圆直径的比较（单位：米）}
\begin{tabular}{lrrc}
\toprule
算例 & $D$ & $2R_*$ & 同直径圆覆盖\\
\midrule
三次测向等边三角形 & 20.000000 & 23.094011 & 否\\
两次测向四边形反例 & 34.514803 & 34.532121 & 否\\
两次测向四边形对照例 & 49.385426 & 49.385426 & 是\\
\bottomrule
\end{tabular}
\end{table}
数值计算采用 $10^{-7}$ 米距离容差，表中结果来自合成测向算例。精确三角形反例与数值核验共同表明：\textbf{以定位区域直径为直径的圆不一定能覆盖该区域；是否能够覆盖，应由式\eqref{eq:cov-check}逐顶点判断。}

\vspace{2pt}
{\footnotesize\noindent [1] Welzl E. Smallest enclosing disks (balls and ellipsoids)[C]. New Results and New Trends in Computer Science. LNCS 555, 1991: 359--370.}
